\chapter{Introduzione}
\label{cap:introduzione}
Il lavoro descritto in questa tesi è stato svolto nell'ambito del tirocinio curriculare presso il centro interdipartimentale J-Tech, all'interno di un'attività di progettazione di un banco per la caratterizzazione di strumenti endodontici. La macchina di base e la cella di carico erano già stati stabiliti all'avvio dell'attività; il contributo qui documentato riguarda la progettazione dei componenti di interfaccia e la loro verifica strutturale e la scelta della tavola di posizionamento.
% Traccia (scommenta le sezioni man mano che le scrivi):

\section{Il trattamento endodontico e gli strumenti rotanti in NiTi}
La scelta di uno strumento endodontico rotante in lega nichel-titanio è oggi
affidata in larga parte all'esperienza del clinico. Il mercato offre un numero
elevato di sistemi, ma non esiste un protocollo
standardizzato che ne consenta il confronto delle proprietà meccaniche su basi
omogenee. Il paragone fra sistemi diversi si fonda perciò su
prove condotte con apparati e condizioni non uniformi, i cui risultati sono
difficilmente trasferibili da uno studio all'altro.

A questa disomogeneità di metodo se ne affianca una di natura sperimentale. Le
prove in vitro impiegano comunemente blocchetti in resina con canale artificiale,
preferiti ai denti estratti per la loro uniformità; anche i blocchetti
commerciali, tuttavia, presentano deviazioni di produzione non trascurabili. 
l'Endo-Training-Bloc-J, assunto come riferimento in questo lavoro, produce blocchetti endodontici tramite stampa 3D e afferma che i coefficienti
di variazione dei parametri geometrici del canale diminuiscano drasticamente grazie a questo metodo fornendo una riproducibilità maggiore durante le prove \cite{smutny2022}.

Il quadro si complica ulteriormente se si considera la natura della
sollecitazione. In esercizio lo strumento è soggetto simultaneamente a flessione
alternata, imposta dalla curvatura del canale, e a torsione, dovuta
all'ingaggio della punta sulle pareti dentinali. Le due sollecitazioni vengono
però quasi sempre indagate separatamente, e la vita utile stimata sulla sola
fatica flessionale ne risulta sovrastimata: a parità di curvatura, il tempo a
rottura diminuisce al crescere della coppia applicata \cite{iacono2021}.

Da queste considerazioni discende l'obiettivo del presente lavoro: la
progettazione di un banco che consenta di sottoporre uno strumento endodontico
rotante a carichi combinati di fatica flessionale e torsione statica, in
condizioni note e ripetibili, misurando la coppia resistente fino a rottura.
% \section{Modalita' di rottura: fatica flessionale e torsione statica}
\section{Obiettivo: un banco a carichi combinati}

\section{Struttura della tesi}
